\subsection{Datasets}\label{sec:data}
\owner{each dataset's owner writes its paragraph (max.\ 120 words)}

% Rubric: "target challenging datasets". Say why each dataset is hard for at least one method.

\begin{table}[h]\centering\small
\caption{Datasets, grouped by the data property each one tests. \decide{final list; drop or keep Mall;
Ky's 1965--2016 earthquakes replaced by the 2023 surface-distance pipeline?}}
\label{tab:data}
\resizebox{\textwidth}{!}{\begin{tabular}{lllrrll}
\toprule
\textbf{ID} & \textbf{Dataset} & \textbf{Owner} & \textbf{$n$} & \textbf{$d$} & \textbf{Labels} & \textbf{Property tested} \\
\midrule
D1 & make\_moons (synthetic)          & Lam            & 1{,}000  & 2        & 2 classes    & curved shapes (control) \\
D2 & Earthquakes, 2023, M$\geq$4.5    & Wei Yew        & 7{,}643  & 2        & none         & long belts, uneven density \\
D3 & Uber NYC pickups, Apr 2014       & Ky, Hoang Anh  & \decide{6k or 15k} & 2 & none      & density contrast in a city \\
D4 & Ecoli                            & Ky             & 336      & 7        & 8 classes    & dimensionality, rare classes \\
D5 & Fashion-MNIST (PCA)              & Lam            & 3{,}000  & 50       & 10 classes   & high dimensions \\
D6 & Wine                             & Lam, Hoang Anh & 178      & 13       & 3 classes    & silhouette vs labels \\
D7 & Breast cancer (Wisconsin)        & Hoang Anh      & 569      & 30       & 2 classes    & silhouette vs labels \\
D8 & HDB resale flats, 2017--2026     & Wei Yew        & 20{,}000\textsuperscript{b} & 4 & flat type (7) & no gaps, scale, recorded ties \\
\bottomrule
\end{tabular}}
\par\raggedright\footnotesize\textsuperscript{b}\,Random sample of 234{,}294 usable sales; all 234{,}294 used for runtime (Section~\ref{sec:runtime}).
\end{table}

\paragraph{D1 make\_moons.} \todo{Lam}{generator settings (noise 0.07), why it is the shape control.}
\paragraph{D2 Earthquakes.}
% Source: weiyew/sc4020-earthquake/notebooks/eda.ipynb, sections 1, 4, 8 and 12.
Every event of magnitude 4.5 or more in 2023 from the USGS ComCat catalogue~\cite{usgs2017}: 7{,}651 events, of which
8 volcanic eruptions near Saipan are removed, leaving 7{,}643 earthquakes. We cluster the epicentre only.
Depth is left out because 44.2\% of events carry the default depth of 10\,km. Each method meets a
different difficulty. Earthquakes lie along narrow belts thousands of kilometres long, which breaks the
compact-cluster assumption of K-means. Density varies widely: the distance to the nearest earthquake
runs from 1.2\,km (5th percentile) to 100.5\,km (95th), so no single DBSCAN radius suits every belt.
And 17.2\% of earthquakes lie within $10^\circ$ of the $180^\circ$ meridian, where flat coordinates
cut the map. There are no labels.
\paragraph{D3 Uber NYC.} \todo{Ky + Hoang Anh}{one merged paragraph; agree sample size and bounding box.}
\paragraph{D4 Ecoli.} \todo{Ky}{Port from Sec.~4.1.}
\paragraph{D5 Fashion-MNIST.} \todo{Lam}{3{,}000-image sample, PCA to 50 dims; why subsampled.}
\paragraph{D6 Wine.} \todo{Lam + Hoang Anh}{one merged paragraph; both runs give K-Means ARI 0.897.}
\paragraph{D7 Breast cancer.} \todo{Hoang Anh}{30 features, standardised; benign vs malignant.}
\paragraph{D8 HDB resale.}
% Source: weiyew/sc4020-hdb/notebooks/eda.ipynb, sections 6, 10, 11 and 17;
% sample size: weiyew/sc4020-hdb/notebooks/clustering.ipynb, section 1;
% floor-area ARI: weiyew/sc4020-hdb/results/screening_labels.csv.
Every resale of a public housing flat in Singapore from January 2017, from data.gov.sg~\cite{hdb2026}. The 6{,}051
sales after June 2026 (2.5\%) have no price index value yet and are dropped, leaving 234{,}294. Methods
are compared on a random sample of 20{,}000 (seed 0); all 234{,}294 are used for runtime. There are four
features: floor area, remaining lease, storey, and price adjusted to 2026 levels with the HDB resale
price index. Flat type (7 classes) is held out as a label, but it is nearly a function of size: floor
area alone reaches ARI 0.601 against it. The data are hard for density methods. There are no gaps
between groups of flats, and storeys are recorded in 3-storey bands, so 77.7\% of sales share their
floor area, lease and storey with at least one other sale.

\subsection{Preprocessing}\label{sec:prep}
\owner{Wei Yew}
Three rules apply wherever a dataset has the property they address. Each is tested by an ablation in
Section~\ref{sec:ablation}.

\paragraph{Coordinates on the sphere.}
% Source: weiyew/sc4020-earthquake/notebooks/eda.ipynb, section 12.
Latitude and longitude treated as flat numbers put two earthquakes 6.07\,km apart at $359.94^\circ$ of
longitude from each other, because they lie on either side of the $180^\circ$ meridian. Density methods
therefore use the haversine (great-circle) distance on coordinates in radians. K-means needs a vector
space to average in, so each epicentre becomes a 3D point on the unit sphere. The straight-line distance
between two such points grows with their great-circle distance, so the clustering follows surface distance
(Ablation~B).

\paragraph{Standardised features.}
% Source: weiyew/sc4020-hdb/notebooks/eda.ipynb, section 9.
Features in their raw units are not comparable. On D8, price (in dollars) accounts for 100.0\% of the
squared distance between sales, so the other three features are ignored. Every feature is scaled to mean
0 and standard deviation 1 (Ablation~C).

\paragraph{Values recorded in steps.}
% Source: weiyew/sc4020-hdb/hdb/features.py, jitter_ties.
When a feature is recorded in steps, many points share exact values, and density methods read each shared
value as a dense spot. Before running a density method on D8, each stepped value is moved by a uniform
random amount of up to half its step: $\pm1.5$ storeys, $\pm0.5$\,m$^2$ of floor area and
$\pm\frac{1}{24}$ year of lease. Price is continuous and is left unchanged. Results are reported with and
without this jitter (Ablation~E).
